\documentclass{article}
\usepackage{graphicx}
\usepackage{amsmath}
\usepackage{amsfonts}
\usepackage{amssymb}
\usepackage{bm}
\usepackage{xcolor}
\usepackage[colorlinks=true,linkcolor=blue!60!black]{hyperref}

\newcommand{\eqname}{Eq.}
\newcommand{\correction}[1]{\textcolor{red}{#1}}
\newcommand{\Nc}{N_c}
\newcommand{\Nf}{N_f}
\newcommand{\Kern}{\mathcal K}
\newcommand{\ASR}{\mathcal S}
\newcommand{\Perm}{\mathcal P}
\newcommand{\Ext}{\mathcal E}

\title{Interpolating the q-mesh in TD-SCHA}
\author{Lorenzo Monacelli}
\date{July 2026}

\begin{document}

\maketitle

\begin{abstract}
The anharmonic phonon spectral function controls experimentally relevant
properties such as thermal transport and Raman, infrared, and neutron
spectra. Its convergence requires dense momentum meshes because the
two-phonon continuum is sampled by pairs of wavevectors. This is especially
difficult in TD-SCHA, where third- and fourth-order vertices enter without
being stored explicitly.

We formulate the interpolation as an image-assignment kernel acting on the
stochastic vectors already used by TD-SCHA. A low-rank factorization of this
geometry kernel becomes a sum of windowed Bloch transforms. The force
constants themselves are never fitted. We identify the linear constraints
that make the interpolated tensor exactly equal to the original tensor at
every wavevector commensurate with the coarse supercell, preserve
permutation symmetry, and preserve the acoustic sum rule (ASR) on every
tensor leg. The construction is given explicitly for tensors of rank two,
three, and four. We then propose a constrained, geometry-only fitting
strategy and discuss its setup cost, runtime cost, and the changes required
in the current TD-SCHA implementation.
\end{abstract}

\section{Introduction}
\label{sec:introduction}

A periodic supercell does not determine which replica of an interaction
cluster should be used away from its commensurate q mesh. Fourier
interpolation therefore needs a rule that breaks the artificial supercell
periodicity. We call this rule the centering kernel.

The central idea is to factorize the centering kernel instead of constructing
a high-order force-constant tensor. Each factor multiplies one stochastic
real-space vector and defines a windowed Bloch transform. The existing
TD-SCHA estimator is multilinear in those vectors, so a sum of separable
kernel terms reproduces the contraction with the corresponding centered
tensor.

Three properties must be imposed on the assembled kernel:
\begin{enumerate}
    \item it must reduce exactly to the original coarse-supercell tensor at
    all commensurate q points;
    \item it must obey the permutation symmetry of the tensor, including the
    corresponding relabeling of momenta and stochastic slots;
    \item it must map every periodic tensor satisfying the ASR to an extended
    tensor satisfying the ASR on every leg.
\end{enumerate}
These are hard linear constraints. Locality, by contrast, is the fitting
objective. Keeping this distinction is essential: fitting a compact
minimal-image target and repairing the ASR afterward generally changes the
kernel that was fitted.

\section{Centering as an image-assignment kernel}
\label{sec:kernel}

\subsection{Lattice and Fourier conventions}
\label{sec:notation}

Let $\bm R$ denote a primitive-cell lattice vector inside the coarse
supercell, and let
\begin{equation}
    \bm L(\bm p)=p_1\bm L_1+p_2\bm L_2+p_3\bm L_3,
    \qquad \bm p\in\mathbb Z^3,
    \label{eq:supercell-translation}
\end{equation}
be a supercell lattice vector. A folded index $i=(a,\alpha,\bm R)$ contains
the primitive atom $a$, Cartesian component $\alpha$, and coarse-cell
coordinate $\bm R$. Its extended replicas are denoted by
$(i,\bm p)$ and have Cartesian position
\begin{equation}
    \bm x(i,\bm p)=\bm\tau_a+\bm R+\bm L(\bm p),
    \label{eq:extended-position}
\end{equation}
where $\bm\tau_a$ is the basis position. All lattice coordinates in phases
are fractional, and every Fourier phase contains $2\pi$ explicitly.

For a real-space stochastic vector $v_i$, the coarse Bloch transform is
\begin{equation}
    w_a^\alpha(\bm q)
    =\frac{1}{\sqrt{\Nc}}\sum_{\bm R}
    e^{-2\pi i\bm q\cdot\bm R}v_a^\alpha(\bm R),
    \label{eq:coarse-bloch}
\end{equation}
where $\Nc$ is the number of primitive cells in the coarse supercell.
The number $\Nf$ of points in the fine q mesh does not enter this transform.
It enters the fine-mesh quadrature and the two-phonon pair normalization,
which are separate from the centering constraints.

\subsection{Rank-two example}
\label{sec:rank2}

Here $i=(a,\alpha,\bm R_i)$ and $j=(b,\beta,\bm R_j)$ are complete
degrees of freedom of the coarse supercell: they specify the atom, Cartesian
component, and primitive-cell position inside the supercell. The integer
vectors $\bm p_1,\bm p_2\in\mathbb Z^3$ select two periodic replicas.
Thus, the infinite-crystal positions associated with the two tensor legs
are $\bm R_i+\bm L(\bm p_1)$ and
$\bm R_j+\bm L(\bm p_2)$, with the corresponding atomic basis positions
understood.

The symbol $\Phi^{(2)}_{ij}$ denotes the single periodic matrix element
stored in the coarse supercell. It has no image arguments because
supercell periodicity identifies all its replicas. In contrast,
$\varPhi^{(2)}_{ij}(\bm p_1,\bm p_2)$ denotes the centered matrix element
assigned to the particular pair of infinite-crystal replicas selected by
$\bm p_1$ and $\bm p_2$. The two objects are related by
\begin{equation}
    \varPhi^{(2)}_{ij}(\bm p_1,\bm p_2)
    =
    \Phi^{(2)}_{ij}
    K^{(2)}_{ij}(\bm p_1,\bm p_2).
    \label{eq:rank2-extension}
\end{equation}
Therefore $K^{(2)}_{ij}(\bm p_1,\bm p_2)$ is the dimensionless fraction of
the periodic element $\Phi^{(2)}_{ij}$ assigned to that replica pair. For
example, a closest-image kernel assigns weight one to the closest pair and
zero to the other replicas; exact geometrical ties share the unit weight.

Only the relative image $\bm p_2-\bm p_1$ is physical, because shifting
both legs by the same supercell translation does not change the pair. One
may fix $\bm p_1=\bm 0$, but it is safer to impose ASR and permutation
constraints in the translation-covariant two-image representation and
remove the common translation afterward.

For a pair, the geometry target assigns unit total weight to the closest
replica or splits that weight equally among exact ties. If
$\mathcal M_{ij}$ is the set of minimizing relative images, then
\begin{equation}
    T^{(2)}_{ij}(\bm p)
    =
    \begin{cases}
    |\mathcal M_{ij}|^{-1}, & \bm p\in\mathcal M_{ij},\\
    0, & \text{otherwise}.
    \end{cases}
    \label{eq:rank2-target}
\end{equation}
This is a target, not yet the final kernel, because centering can conflict
with the ASR at finite support.

Flattening $(i,\bm p_1)$ and $(j,\bm p_2)$ into row and column indices makes
$K^{(2)}$ a matrix. It therefore admits the factorization
\begin{equation}
    K^{(2)}_{ij}(\bm p_1,\bm p_2)
    =
    \sum_{\xi=1}^{N_r}
    c_\xi A_i^{(\xi)}(\bm p_1)
    B_j^{(\xi)}(\bm p_2).
    \label{eq:rank2-factorization}
\end{equation}
For rank two this representation is exact when $N_r$ reaches the matrix
rank. This matrix-rank statement does not extend unchanged to rank-three
and rank-four canonical polyadic decompositions.

The factors define windowed Bloch vectors
\begin{align}
    w_a^{A\xi,\alpha}(\bm q)
    &=
    \frac{1}{\sqrt{\Nc}}
    \sum_{\bm R,\bm p}
    A_{a\alpha\bm R}^{(\xi)}(\bm p)
    e^{-2\pi i\bm q\cdot[\bm R+\bm L(\bm p)]}
    v_a^\alpha(\bm R),                                      \label{eq:window-a}\\
    w_b^{B\xi,\beta}(-\bm q)
    &=
    \frac{1}{\sqrt{\Nc}}
    \sum_{\bm R,\bm p}
    B_{b\beta\bm R}^{(\xi)}(\bm p)
    e^{+2\pi i\bm q\cdot[\bm R+\bm L(\bm p)]}
    v_b^\beta(\bm R).                                      \label{eq:window-b}
\end{align}
The replica-resolved element in
\autoref{eq:rank2-extension} is a real-space quantity. Its q-space Fourier
transform is denoted by $\varPhi^{(2)}_{ab}(\bm q)$ below. The cell and image
coordinates have been Fourier summed, and the Cartesian components are
suppressed to keep the expression readable. The superscript $(2)$ denotes
the tensor order, not a power. With these conventions,
\begin{equation}
    \varPhi^{(2)}_{ab}(\bm q)
    =
    \sum_{\xi=1}^{N_r}c_\xi
    \left\langle
    w_a^{(1)A\xi}(\bm q)
    w_b^{(2)B\xi}(-\bm q)
    \right\rangle .
    \label{eq:phi2-interp}
\end{equation}
The opposite signs in \autoref{eq:window-a} and
\autoref{eq:window-b} are required by momentum conservation and by the
Hermitian relation between $\bm q$ and $-\bm q$.

\subsection{Corrections to the preliminary derivation}
\label{sec:corrections}

Several distinctions are needed before generalizing the construction.

First, the centering kernel depends on atom, Cartesian, folded-cell, and
image labels. A kernel written only as $K_{ab}(p,q,r)$ does not contain
enough information to describe atom-resolved ties or a general interaction
cluster.

Second, the translational average is over the $\Nc$ coarse cells. The
quantity $\Nf$ counts fine q points and must not replace that average.
Factors of $\Nf/\Nc$ associated with the TD-SCHA vertex normalization are
applied once in the estimator; they are not a window-normalization
condition.

Third, a constant image sum for every individual factor is neither the
general ASR condition nor a necessary one. ASR, commensurate identity, and
permutation symmetry constrain the assembled kernel
$\sum_\xi c_\xi A^{(\xi)}\otimes B^{(\xi)}$, and different low-rank terms
may cancel in satisfying those constraints.

Fourth, adding permutations of factors is correct only if the atom,
Cartesian, image, momentum, and stochastic-slot labels are permuted
together and the orbit has the correct multiplicity. This becomes important
for the three force channels of the rank-three estimator and the four force
channels of the rank-four estimator.

\section{Exact constraints}
\label{sec:constraints}

\subsection{Identity at commensurate q points}
\label{sec:commensurate}

At a q point commensurate with the coarse supercell,
\begin{equation}
    e^{2\pi i\bm q\cdot\bm L(\bm p)}=1
    \quad\text{for every }\bm p.
    \label{eq:commensurate-phase}
\end{equation}
The extended tensor is therefore exactly equal to the periodic tensor for
arbitrary tensor entries if and only if the image weights form a partition
of unity. For rank $n$ this condition is
\begin{equation}
    \sum_{\bm p_1,\ldots,\bm p_n}
    K^{(n)}_{i_1\ldots i_n}
    (\bm p_1,\ldots,\bm p_n)
    =1
    \quad\text{for every }(i_1,\ldots,i_n).
    \label{eq:partition}
\end{equation}
After fixing a common translation, the redundant common-image sum is
removed and the same condition applies to the $n-1$ relative images.

Equation~\eqref{eq:partition} is stronger and clearer than matching a
normalization only after an ensemble average. It proves equality
configuration by configuration and for every periodic tensor. It also
shows why the condition contains no $\Nf/\Nc$ factor.

\subsection{Permutation symmetry}
\label{sec:permutation}

Let $\pi$ be a permutation of the $n$ legs. The kernel must satisfy
\begin{equation}
    K^{(n)}_{i_{\pi(1)}\ldots i_{\pi(n)}}
    (\bm p_{\pi(1)},\ldots,\bm p_{\pi(n)})
    =
    K^{(n)}_{i_1\ldots i_n}
    (\bm p_1,\ldots,\bm p_n).
    \label{eq:permutation}
\end{equation}
For a separable term, this is enforced by its permutation orbit
\begin{equation}
    \operatorname{Sym}
    \left[
    A_1\otimes\cdots\otimes A_n
    \right]
    =
    \frac{1}{n!}\sum_{\pi\in S_n}
    A_{\pi(1)}\otimes\cdots\otimes A_{\pi(n)},
    \label{eq:symmetrized-term}
\end{equation}
where every physical and momentum label follows the same permutation.
Repeated factors reduce the number of distinct runtime terms.

\subsection{Can the same centering factor be used on every leg?}
\label{sec:symmetric-powers}

Yes, permutation symmetry suggests the more restrictive decomposition
\begin{equation}
    K^{(n)}_{i_1\ldots i_n}
    (\bm p_1,\ldots,\bm p_n)
    \simeq
    \sum_{\xi=1}^{R_n^{\mathrm{sym}}}c_\xi
    \prod_{s=1}^{n}
    A_{i_s}^{(\xi)}(\bm p_s).
    \label{eq:symmetric-power-kernel}
\end{equation}
Within one term, the same centering factor $A^{(\xi)}$ is evaluated on
every leg. Exchanging any two complete leg labels
$(i_s,\bm p_s)$ leaves the product unchanged, so
\autoref{eq:symmetric-power-kernel} satisfies permutation symmetry
automatically. No explicit sum over $n!$ permutations is needed.

Here ``the same factor'' does not mean that the stochastic data are the
same. A force leg still contains a force field and a displacement leg still
contains a displacement field. They are merely multiplied by the same
geometry window $A^{(\xi)}$ before entering the multilinear estimator.
The momenta on different legs also remain different.

For rank two the result is exact and elementary. Combine $(i,\bm p)$ into
one matrix index. A real permutation-invariant kernel is a real symmetric
matrix, so its spectral decomposition is
\begin{equation}
    K^{(2)}
    =
    Q\Lambda Q^T
    =
    \sum_{\xi=1}^{R_2}
    \lambda_\xi q_\xi q_\xi^T.
    \label{eq:symmetric-rank2-decomposition}
\end{equation}
Thus one may set $A^{(\xi)}=B^{(\xi)}=q_\xi$ and
$c_\xi=\lambda_\xi$ in \autoref{eq:rank2-factorization} without losing any
degree of freedom allowed by symmetry. The coefficients must be allowed to
have either sign. A kernel can have non-negative image weights and still be
an indefinite matrix.

For ranks three and four, sums of equal-factor terms are symmetric-power
decompositions. They remain complete in the following sense: every
finite-dimensional fully symmetric tensor can be written as a finite sum
of powers $A_\xi^{\otimes n}$. One way to see this is the polarization
identity
\begin{align}
    &\operatorname{Sym}
    [A_1\otimes\cdots\otimes A_n]
    \nonumber\\
    &\quad =
    \frac{1}{2^n\,n!}
    \sum_{\boldsymbol\epsilon\in\{-1,+1\}^n}
    \left(\prod_{s=1}^n\epsilon_s\right)
    \left(\sum_{s=1}^n\epsilon_s A_s\right)^{\otimes n}.
    \label{eq:polarization}
\end{align}
Therefore any symmetrized product of different factors can be replaced by
a sum of terms that use the same factor on every leg.

Completeness does not imply equal algebraic rank. One symmetrized
term with distinct factors may require as many as $2^{n-1}$ powers
after pairing the $\boldsymbol\epsilon$ and $-\boldsymbol\epsilon$
contributions in \autoref{eq:polarization}. This gives at most four powers
for rank three and eight for rank four. The minimum symmetric rank need not
equal the minimum unrestricted canonical-polyadic rank.

The algebraic rank is not, however, the runtime cost. A distinct-factor
term represents
\begin{equation}
    \frac{1}{n!}\sum_{\pi\in S_n}
    A_{\pi(1)}\otimes\cdots\otimes A_{\pi(n)}.
    \label{eq:explicit-permutation-orbit}
\end{equation}
If this sum is applied literally, it needs up to $n!$ ordered window
products. An equal-factor power is already invariant and needs one ordered
product, not another permutation sum. Consequently the polarization
identity replaces, in the worst generic case,
\begin{equation}
    6\ \text{ordered products by at most }4
    \quad(n=3),
    \qquad
    24\ \text{by at most }8
    \quad(n=4).
    \label{eq:polarization-runtime-bound}
\end{equation}
Thus the factor of four or eight is an upper bound on the number of
equal-factor terms needed to replace one general symmetrized term; it is not
a claim that the fit uses four or eight times fewer algebraic terms. When
permutations are evaluated separately, the equal-factor representation can
still use fewer runtime products because it avoids the $n!$ orbit.

This bound is not a performance guarantee. Repeated factors reduce the
number of distinct permutations, and a specialized kernel can reuse
intermediates or accumulate several permutations in one call. Independently
optimized equal-factor and distinct-factor fits can also converge at
different ranks. The quantities that must be compared are therefore
\begin{equation}
    R_{n,\mathrm{sym}}^{\mathrm{run}}
    =
    R_n^{\mathrm{sym}},
    \qquad
    R_{n,\mathrm{gen}}^{\mathrm{run}}
    =
    \sum_{\xi=1}^{R_n^{\mathrm{gen}}}P_\xi,
    \qquad 1\leq P_\xi\leq n!,
    \label{eq:runtime-term-counts}
\end{equation}
where $P_\xi$ is the number of distinct ordered products evaluated
for general term $\xi$ after merging equal permutations and accounting for
kernel reuse. These measured runtime counts, together with field-generation
and contraction time, decide which representation is advantageous.

For even order, including rank four, signed coefficients are essential. If
$c_\xi\geq0$ were imposed, then
\begin{equation}
    K^{(n)}(u,\ldots,u)
    =
    \sum_\xi c_\xi
    [A^{(\xi)}(u)]^n
    \geq0
    \qquad\text{for even }n,
    \label{eq:positive-power-restriction}
\end{equation}
which restricts the fit to positive semidefinite forms. The projected
centering kernel is not guaranteed to have this property.

The equal-factor form imposes permutation symmetry but does not by itself
impose translation covariance, partition of unity, or ASR. Those remain
the hard linear constraints described below. Symmetry does reduce their
number: after exact permutation symmetry is built in, enforcing an ASR
condition on one leg implies its permuted counterparts on the other legs.

This ansatz is especially attractive for D4. A rank-four term of the form
in \autoref{eq:symmetric-power-kernel} can be evaluated with the current
reused $w,v$ field path because all four logical legs use the same window.
A four-slot Julia interface is required only if the constrained
symmetric-power rank is too large and the more efficient distinct-factor
decomposition is needed. The geometry fit should therefore compare these
two representations before the runtime interface is expanded.


\subsection{Acoustic sum rule}
\label{sec:asr}

We first write the ASR of the periodic tensor without operator notation.
Recall that a coarse-supercell index is
$i_s=(a_s,\alpha_s,\bm R_s)$. Choose a tensor leg $s$ and a Cartesian
translation direction $\beta$. Hold all other tensor indices fixed and set
the Cartesian component on leg $s$ to $\beta$. The periodic ASR is
\begin{align}
    &\left(\ASR_s^{\mathrm{per}}\Phi^{(n)}\right)^\beta_
    {i_1\ldots\widehat{i_s}\ldots i_n}
    \nonumber\\
    &\quad\equiv
    \sum_{\bm R_s}
    \sum_{a_s}
    \Phi^{(n)}_
    {i_1\ldots(a_s,\beta,\bm R_s)\ldots i_n}
    =0.
    \label{eq:periodic-asr-sum}
\end{align}
The sum runs over every primitive-cell position $\bm R_s$ and every atom
$a_s$ in the coarse supercell. The hat means that leg $s$ is absent from
the list of free indices. Equation~\eqref{eq:periodic-asr-sum} says that a
rigid translation of all atoms in direction $\beta$ produces no force.

The centered tensor has one additional image label $\bm p_t$ on every leg.
Write $\bar i_t=(i_t,\bm p_t)$. Its ASR is the same physical sum, but now
the translated leg must also be summed over all of its supercell replicas:
\begin{align}
    &\left(\ASR_s^{\mathrm{ext}}\varPhi^{(n)}\right)^\beta_
    {\bar i_1\ldots\widehat{\bar i_s}\ldots\bar i_n}
    \nonumber\\
    &\quad\equiv
    \sum_{\bm p_s\in\mathbb Z^3}
    \sum_{\bm R_s}
    \sum_{a_s}
    \varPhi^{(n)}_
    {\bar i_1\ldots((a_s,\beta,\bm R_s),\bm p_s)\ldots\bar i_n}
    =0.
    \label{eq:extended-asr-sum}
\end{align}
All other folded indices and image labels are held fixed. Although the image
sum is written over $\mathbb Z^3$, only the finite support of the kernel
contributes.

The extension by the centering kernel is not a new physical operation. It
is simply the definition
\begin{align}
    &\left(\Ext_K\Phi^{(n)}\right)_
    {(i_1,\bm p_1)\ldots(i_n,\bm p_n)}
    \nonumber\\
    &\quad\equiv
    \Phi^{(n)}_{i_1\ldots i_n}
    K^{(n)}_{i_1\ldots i_n}
    (\bm p_1,\ldots,\bm p_n)
    =
    \varPhi^{(n)}_
    {(i_1,\bm p_1)\ldots(i_n,\bm p_n)}.
    \label{eq:kernel-extension-explicit}
\end{align}
Substituting \autoref{eq:kernel-extension-explicit} into
\autoref{eq:extended-asr-sum} gives
\begin{align}
    &\left(\ASR_s^{\mathrm{ext}}\Ext_K\Phi^{(n)}\right)^\beta
    \nonumber\\
    &\quad=
    \sum_{\bm R_s}
    \sum_{a_s}
    \Phi^{(n)}_
    {i_1\ldots(a_s,\beta,\bm R_s)\ldots i_n}
    H_s(a_s,\bm R_s),
    \label{eq:extended-asr-expanded}
\end{align}
where the kernel image sum is
\begin{align}
    H_s(a_s,\bm R_s)
    &\equiv
    \sum_{\bm p_s\in\mathbb Z^3}
    K^{(n)}_
    {i_1\ldots(a_s,\beta,\bm R_s)\ldots i_n}
    (\bm p_1,\ldots,\bm p_n).
    \label{eq:asr-image-sum}
\end{align}
The arguments on every other leg are fixed in
\autoref{eq:asr-image-sum}. Therefore $H_s$ may depend on those other
indices and images, on $s$, and on $\beta$. The ASR requirement is only
that it must not depend on the atom $a_s$ or folded cell $\bm R_s$ being
summed on leg $s$:
\begin{equation}
    H_s(a_s,\bm R_s)=\gamma_s
    \quad\text{for every }a_s,\bm R_s.
    \label{eq:asr-constant-image-sum}
\end{equation}
Under this condition, \autoref{eq:extended-asr-expanded} becomes
\begin{align}
    \left(\ASR_s^{\mathrm{ext}}\Ext_K\Phi^{(n)}\right)^\beta
    &=
    \gamma_s
    \sum_{\bm R_s}\sum_{a_s}
    \Phi^{(n)}_
    {i_1\ldots(a_s,\beta,\bm R_s)\ldots i_n}
    \nonumber\\
    &=
    \gamma_s
    \left(\ASR_s^{\mathrm{per}}\Phi^{(n)}\right)^\beta
    =0.
    \label{eq:asr-preservation-explicit}
\end{align}
This is the complete ASR-preservation argument. No additional map is
needed.

In the translation-covariant representation above, the same image sum
applies to every leg. If the common translation is removed by fixing
$\bm p_1=\bm0$, the formulas take two forms. For a non-reference leg
$s>1$, \autoref{eq:asr-image-sum} remains an ordinary sum over
$\bm p_s$. For the reference leg, holding the absolute positions of all
other legs fixed while changing its image by $\bm t$ shifts every relative
image. Its kernel sum is therefore
\begin{align}
    H_1(a_1,\bm R_1)
    &\equiv
    \sum_{\bm t\in\mathbb Z^3}
    K^{(n)}_{(a_1,\beta,\bm R_1)i_2\ldots i_n}
    (\bm0,\bm p_2-\bm t,\ldots,\bm p_n-\bm t).
    \label{eq:asr-reference-diagonal-sum}
\end{align}
This is why the reference-leg ASR appears as a diagonal common shift when
the first image is pinned.

The ASR constraints used in the fit are now explicit. Choose any reference
atom-cell pair $(a_0,\bm R_0)$. For every leg, Cartesian direction, and
choice of the other fixed arguments, add the homogeneous row
\begin{equation}
    H_s(a_s,\bm R_s)-H_s(a_0,\bm R_0)=0.
    \label{eq:asr-fit-row}
\end{equation}
Each $H_s$ is a linear sum of kernel entries, so
\autoref{eq:asr-fit-row}, partition of unity, and permutation symmetry are
all linear equations. After vectorizing the kernel they can be collected as
\begin{equation}
    C_n k_n=d_n,
    \qquad k_n=\operatorname{vec}K^{(n)}.
    \label{eq:linear-constraints}
\end{equation}
Redundant rows are expected because permutation and ASR constraints
overlap. They should be handled with a rank-revealing QR decomposition or
a pseudoinverse, not removed by hand.

\section{Geometry targets for rank three and rank four}
\label{sec:high-rank}

\subsection{A cluster target rather than a force-star target}
\label{sec:cluster-target}

For a folded $n$-tuple, define the positions
$\bm x_s=\bm x(i_s,\bm p_s)$ and the complete-graph cluster cost
\begin{equation}
    D_n(\bm p_1,\ldots,\bm p_n)
    =
    \sum_{1\leq s<t\leq n}
    \omega_{st}\,
    \left|\bm x_s-\bm x_t\right|^\rho ,
    \qquad \rho=1\text{ or }2.
    \label{eq:cluster-cost}
\end{equation}
The metric and weights must be fixed for a calculation. The choice
$\rho=1$ gives a perimeter, while $\rho=2$ gives a smooth squared-distance
cost. Basis positions $\bm\tau_a$ must be included. Cell-index geometry
alone cannot resolve basis-dependent ties in small real-material
supercells.

Let $\mathcal M_n(i_1,\ldots,i_n)$ contain all image tuples minimizing
\autoref{eq:cluster-cost}, modulo a common translation. The raw target is
\begin{equation}
    T^{(n)}_{i_1\ldots i_n}
    (\bm p_1,\ldots,\bm p_n)
    =
    \begin{cases}
    |\mathcal M_n|^{-1},
    &(\bm p_1,\ldots,\bm p_n)\in\mathcal M_n,\\
    0,&\text{otherwise}.
    \end{cases}
    \label{eq:raw-target}
\end{equation}
Exact ties are therefore split without breaking space-group or permutation
symmetry.

Independent minimal images relative to one pinned force leg define a star,
not the complete cluster in \autoref{eq:cluster-cost}. The repository tests
show that this distinction is decisive for rank four: on the SnTe
$2\times2\times2$ geometry, the tested force-star assignment differs from
the four-leg complete-graph target in most atom-cell classes and can move
the fourth-order spectral shift in the wrong direction. The D4 fit must
therefore use four logical legs from the beginning.

\subsection{Rank-three tensor}
\label{sec:rank3}

The same distinction applies at rank three. The indices $i,j,k$ are three
complete coarse-supercell degrees of freedom, while
$\bm p_1,\bm p_2,\bm p_3$ select one replica for each leg. The periodic
entry $\Phi^{(3)}_{ijk}$ is stored only once. If the centered tensor were
formed explicitly, its element would be
\begin{equation}
    \varPhi^{(3)}_{ijk}(\bm p_1,\bm p_2,\bm p_3)
    =
    \Phi^{(3)}_{ijk}
    K^{(3)}_{ijk}(\bm p_1,\bm p_2,\bm p_3).
    \label{eq:rank3-extension}
\end{equation}
Thus $K^{(3)}$ distributes one periodic third-order entry among triplets of
infinite-crystal replicas. The tensor-free algorithm never forms
$\varPhi^{(3)}$; it realizes the same multiplication through three
windowed stochastic fields.

For rank three, the target is the minimal triangle with edges
$(1,2)$, $(1,3)$, and $(2,3)$. Since the target and constraints are already
permutation symmetric, the primary representation is
\begin{equation}
    K^{(3)}_{ijk}
    (\bm p_1,\bm p_2,\bm p_3)
    \simeq
    \sum_{\xi=1}^{R_3}c_\xi
    A_i^{(\xi)}(\bm p_1)
    A_j^{(\xi)}(\bm p_2)
    A_k^{(\xi)}(\bm p_3).
    \label{eq:rank3-factorization}
\end{equation}
Each rank term is invariant under exchange of complete leg labels, so no
six-term permutation orbit is applied at runtime. The three stochastic
force channels still place the force field on the $z$, $w$, or $v$ leg,
but all three use the same geometry factor $A^{(\xi)}$.

After Fourier transforming the replica coordinates, the following equation
uses $\varPhi^{(3)}$ for the centered q-space tensor. Its three displayed
arguments are the momenta carried by its legs. Atom and Cartesian indices
are suppressed because TD-SCHA contracts them with mode eigenvectors. The
windowed estimator has the form
\begin{equation}
    \varPhi^{(3)}(\bm q_1,\bm q_2,\bm q_3)
    =
    \sum_{\xi=1}^{R_3}c_\xi
    \left\langle
    w^{A\xi}(\bm q_1)
    w^{A\xi}(\bm q_2)
    w^{A\xi}(\bm q_3)
    \right\rangle,
    \qquad
    \bm q_1+\bm q_2+\bm q_3=\bm G.
    \label{eq:rank3-estimator}
\end{equation}
The stochastic field type on each leg is selected by the force channel; the
centering factor is the same. Equation~\eqref{eq:rank3-estimator} is
therefore permutation symmetric term by term.

\subsection{Rank-four tensor}
\label{sec:rank4}

At rank four, $i,j,k,l$ denote four complete coarse-supercell degrees of
freedom and $\bm p_1,\ldots,\bm p_4$ select their replicas. The periodic
entry $\Phi^{(4)}_{ijkl}$ and the corresponding centered element are related
by
\begin{equation}
    \varPhi^{(4)}_{ijkl}
    (\bm p_1,\bm p_2,\bm p_3,\bm p_4)
    =
    \Phi^{(4)}_{ijkl}
    K^{(4)}_{ijkl}
    (\bm p_1,\bm p_2,\bm p_3,\bm p_4).
    \label{eq:rank4-extension}
\end{equation}
Hence $K^{(4)}$ assigns the periodic fourth-order entry to a four-replica
cluster. As at rank three, neither $\Phi^{(4)}$ nor $\varPhi^{(4)}$ is
constructed by the tensor-free implementation; only the geometry kernel is
factorized and applied to the four stochastic fields.

For rank four, the target contains all six edges of the four-vertex
cluster. The four logical legs in the TD-SCHA two-phonon operation are
\begin{equation}
    (E_w,E_v,I_w,I_v),
    \label{eq:d4-legs}
\end{equation}
where $E_w,E_v$ are the external two-phonon legs and $I_w,I_v$ are the
internal legs contracted with the input two-phonon vector. Their momenta
can be written as
\begin{equation}
    (\bm q_1,\bm q_2,-\bm q_a,-\bm q_b),
    \qquad
    \bm q_1+\bm q_2
    =
    \bm q_a+\bm q_b
    =
    \bm q_{\text{pert}}+\bm G.
    \label{eq:d4-momenta}
\end{equation}

Because the target is already permutation symmetric, the production
factorization is
\begin{align}
    &K^{(4)}(E_w,E_v,I_w,I_v) \nonumber\\
    &\qquad\simeq
    \sum_{\xi=1}^{R_4}c_\xi
    A^{(\xi)}(E_w)
    A^{(\xi)}(E_v)
    A^{(\xi)}(I_w)
    A^{(\xi)}(I_v).
    \label{eq:rank4-factorization}
\end{align}
Every rank term is invariant under all 24 leg permutations before it is
applied. The external and internal fields have different momenta and
different roles, but they use the same geometry window.

The current implementation can evaluate
\autoref{eq:rank4-factorization} through its reused $w,v$ field path. The
four existing force-position gates place the force field on $E_w,E_v,I_w$,
or $I_v$ while leaving the window unchanged. No four-slot field interface
is needed for this primary representation.

Only if the required symmetric rank is too large should the fallback
\begin{equation}
    K^{(4)}_{\mathrm{gen}}
    =
    \sum_\xi c_\xi
    \operatorname{Sym}
    [A_\xi\otimes B_\xi\otimes C_\xi\otimes D_\xi]
    \label{eq:rank4-distinct-fallback}
\end{equation}
be considered. That form may reduce algebraic rank but requires distinct
logical field inputs and up to 24 ordered products per unsymmetrized term.
Its implementation is justified only by a measured runtime advantage at
the same constrained kernel error.

This factorization is of the geometry kernel, not of $\Phi^{(4)}$. Gaussian
integration by parts relates the ensemble average of the four stochastic
fields to the fourth-order vertex in the infinite-sampling SSCHA limit.
Multilinearity then implies that applying
\autoref{eq:rank4-factorization} to the fields is equivalent to contracting
the fourth-order vertex with the same image kernel. No entry of
$\Phi^{(4)}$ is needed during fitting or application.

\section{Constrained fitting strategy}
\label{sec:fitting}

\subsection{Project the geometry target onto the physical affine space}
\label{sec:projection}

The raw minimal-cluster target generally violates ASR because centering and
ASR do not commute. The correct target for factorization is the closest
kernel satisfying all hard constraints:
\begin{equation}
    k_n^\star
    =
    \underset{k}{\operatorname{argmin}}\;
    \frac{1}{2}(k-t_n)^T W_n(k-t_n)
    \quad\text{subject to}\quad
    C_nk=d_n,
    \label{eq:affine-projection}
\end{equation}
where $t_n=\operatorname{vec}T^{(n)}$ and $W_n$ is a positive metric.
A distance-decay metric may be used to control where the ASR correction is
placed, but it must not change the constraint rows.

For a manageable geometry the projected solution is
\begin{equation}
    k_n^\star
    =
    t_n-W_n^{-1}C_n^T
    \left(C_nW_n^{-1}C_n^T\right)^+
    (C_nt_n-d_n),
    \label{eq:projection-solution}
\end{equation}
where $+$ denotes the pseudoinverse. Before fitting, the code must report
the residuals of partition, ASR, and permutation constraints and the
distance $\|k_n^\star-t_n\|_{W_n}$. A large projection distance means that
the chosen support or geometry target is incompatible with the physical
constraints and should not be hidden by increasing the factor rank.

\subsection{Fit the assembled kernel, not each factor separately}
\label{sec:constrained-cp}

The physical purpose of the fit is simple. The plain kernel already
preserves the coarse-grid tensor and the ASR, but has poor off-grid
locality. We want to add a linear combination of localized kernel patterns
that moves it toward the projected geometry target without changing any of
those exact properties.

We first define every object used in the fit. Let $N_K$ be the number of
kernel entries after choosing the atom, cell, and image support, and let
$N_C$ be the number of linear constraint equations. Then:
\begin{itemize}
    \item $k_n^\star\in\mathbb R^{N_K}$ is the vectorized, ASR-projected
    geometry target obtained from \autoref{eq:affine-projection};
    \item $k_0\in\mathbb R^{N_K}$ is the vectorized plain kernel;
    \item $u_\xi\in\mathbb R^{N_K}$ is one candidate correction pattern in
    the space of kernels;
    \item $R_n$ is the number of candidate correction patterns;
    \item $U=(u_1,\ldots,u_{R_n})\in\mathbb R^{N_K\times R_n}$ is the
    matrix whose columns are those patterns;
    \item $c=(c_1,\ldots,c_{R_n})^T\in\mathbb R^{R_n}$ contains their
    scalar amplitudes;
    \item $C_n\in\mathbb R^{N_C\times N_K}$ and
    $d_n\in\mathbb R^{N_C}$ contain the partition, ASR, and permutation
    equations defined in \autoref{eq:linear-constraints}.
\end{itemize}
The vector $u_\xi$ is not a displacement or force. It is the complete
image-assignment kernel produced by one proposed set of one-leg window
factors, flattened into a column.

There are two ways to construct a candidate column. The equal-factor
representation uses
\begin{equation}
    u_\xi^{\mathrm{sym}}
    =
    \operatorname{vec}
    [(A^{(\xi)})^{\otimes n}],
    \label{eq:symmetric-candidate-columns}
\end{equation}
which is already invariant under leg permutations. The distinct-factor
representation uses
\begin{equation}
    u_\xi^{\mathrm{gen}}
    =
    \operatorname{vec}
    \operatorname{Sym}
    [A_1^{(\xi)}\otimes\cdots\otimes A_n^{(\xi)}].
    \label{eq:candidate-columns}
\end{equation}
Here $A_s^{(\xi)}$ is the one-leg centering window assigned to leg $s$ of
candidate $\xi$. In the equations below, $U$ may be built from either type
of column.

The fitted kernel is
\begin{equation}
    k_{n,R}=k_0+Uc
    =k_0+\sum_{\xi=1}^{R_n}c_\xi u_\xi.
    \label{eq:baseline-correction}
\end{equation}
Thus $Uc$ is the physical correction to the plain image assignment. A
positive or negative coefficient changes how much of the periodic tensor is
assigned to the clusters represented by $u_\xi$.

Both the target and the fitted kernel must satisfy
$C_nk=d_n$. The plain kernel is chosen to be feasible, so
\begin{equation}
    C_nk_0=d_n.
    \label{eq:plain-kernel-feasible}
\end{equation}
Substituting \autoref{eq:baseline-correction} into the same constraint gives
\begin{align}
    C_n(k_0+Uc)&=d_n,\nonumber\\
    C_nk_0+C_nUc&=d_n,\nonumber\\
    C_nUc&=0.
    \label{eq:homogeneous-constraint}
\end{align}
Equation~\eqref{eq:homogeneous-constraint} therefore has a direct physical
meaning: the correction may redistribute weight off the coarse mesh, but
its total change to every commensurate-identity, ASR, and permutation
constraint must be zero.

The matrix $C_nU\in\mathbb R^{N_C\times R_n}$ measures how the candidate
terms violate the constraints. Its entry $(r,\xi)$ is the change in
constraint $r$ produced by adding candidate $u_\xi$ with unit coefficient.
An individual candidate need not be feasible. Different candidates may
cancel one another. For example, if two terms change one partition sum by
$+1$ and $-1$, equal coefficients preserve that partition sum.

To parameterize all such cancellations, let
$r_0=\dim\operatorname{null}(C_nU)$ and let
\begin{equation}
    Z\in\mathbb R^{R_n\times r_0},
    \qquad C_nUZ=0,
    \label{eq:nullspace-basis}
\end{equation}
contain a basis of the nullspace of $C_nU$. Every allowed coefficient
vector can then be written as
\begin{equation}
    c=Zz,
    \qquad z\in\mathbb R^{r_0}.
    \label{eq:nullspace-coordinates}
\end{equation}
The vector $z$ has no additional physical meaning. It is only the set of
free coordinates used to describe combinations of candidate kernels that
leave all exact constraints unchanged. If $r_0=0$, then $c=0$ is the only
allowed combination: the current candidate set cannot improve on the plain
kernel without breaking a constraint.

For fixed one-leg factors, the remaining fit chooses $z$ to approach the
projected geometry target:
\begin{equation}
    \underset{z\in\mathbb R^{r_0}}{\operatorname{minimize}}\;
    \|k_0+UZz-k_n^\star\|_{W_n}^2.
    \label{eq:variable-projection}
\end{equation}
$W_n$ is the positive weighting matrix introduced in
\autoref{eq:affine-projection}; it defines which kernel entries are most
important in the locality error. Because $C_nUZ=0$, every trial value of
$z$ in \autoref{eq:variable-projection} preserves the commensurate tensor,
ASR, and permutation constraints exactly up to numerical precision.

The complete optimization also changes the one-leg factors
$A_s^{(\xi)}$. After each factor update, the code rebuilds $U$, recomputes
the nullspace basis $Z$, and solves for $z$. This outer problem is
non-convex, so several geometry-informed or random initializations are
needed. Increasing $R_n$ adds candidate directions; it is required when
the nullspace is empty or when the feasible fit remains above the target
error.

The acceptance tests are:
\begin{enumerate}
    \item $\|C_nk_{n,R}-d_n\|_\infty$ is at numerical precision;
    \item the direct permutation residual is at numerical precision;
    \item
    $\|k_{n,R}-k_n^\star\|_{W_n}/\|k_n^\star\|_{W_n}$ is below the
    stochastic-noise target;
    \item increasing the fitted rank does not change the physical response
    beyond its sampling uncertainty.
\end{enumerate}
A low kernel-fit error is not sufficient if either of the first two tests
fails.

\subsection{Matrix-free fitting for large primitive cells}
\label{sec:matrix-free}

The dense objects used in the derivation are not acceptable production
data structures. Let $M$ be the number of atom, Cartesian, folded-cell, and
image states available to one kernel leg. A dense fourth-rank kernel has
$M^4$ entries. Even $M=10^3$ gives $10^{12}$ entries, or about
$8$~TB in double precision, before work arrays. The target
$T^{(4)}$, projected kernel $K^{(4)\star}$, candidate matrix $U$, and
constraint matrix $C_4$ must therefore remain implicit.

The fine-mesh size $\Nf$ does not change $M$ or the real-space kernel size.
It creates a different scaling problem: the pair list and each transformed
field contain $O(\Nf)$ entries. Large primitive cells increase both
$M$ and the band count $n_b=3n_{\mathrm{at}}$, while dense fine meshes
increase the field and pair storage. The implementation must control these
two sources of cost separately.

For small validation cells, explicitly building $T^{(4)}$, $C_4$, and
$K^{(4)\star}$ remains useful because every constraint and geometry class
can be inspected. The production symmetric-power fit instead minimizes
the distance to the raw target directly while retaining the feasible form
\begin{equation}
    k_{n,R}=k_0+UZz,
    \qquad C_nUZ=0.
    \label{eq:matrixfree-feasible-form}
\end{equation}
Because every trial kernel is feasible, minimizing
$\|k_{n,R}-t_n\|_{W_n}$ is equivalent to comparing with the affine
projection of $t_n$ inside the chosen low-rank model. The dense vector
$k_n^\star$ in \autoref{eq:projection-solution} is not needed.

The required operations can be evaluated from the one-leg factors:
\begin{enumerate}
    \item \textbf{Kernel entries.} Given one folded atom tuple and image
    tuple, evaluate
    $K_{n,R}=K_0+\sum_\xi c_\xi\prod_s A^{(\xi)}(i_s,\bm p_s)$ on demand.
    \item \textbf{Geometry target entries.} A geometry oracle receives the
    folded tuple, computes the complete-graph distances, returns the
    minimizing image tuple or tied tuples, and never stores neighboring
    target entries.
    \item \textbf{Factor Gram matrix.} For $W_n=I$, two symmetric powers
    obey
    \begin{equation}
        \left\langle
        (A^{(\xi)})^{\otimes n},
        (A^{(\zeta)})^{\otimes n}
        \right\rangle
        =
        \left\langle A^{(\xi)},A^{(\zeta)}\right\rangle^n.
        \label{eq:symmetric-power-gram}
    \end{equation}
    A separable weight metric gives the same reduction with a weighted
    one-leg inner product. Thus $U^TW_nU$ is built from an
    $R_n\times R_n$ one-leg Gram matrix.
    \item \textbf{Target contractions.}\par
    Stream folded tuples, reduced by symmetry, through the geometry oracle. For each tuple, accumulate the
    target-factor overlap and its factor gradients. Neighbor lists and a
    distance-weighted metric
    prioritize compact clusters. If the complete tuple set is still too
    large, reproducible stratified sampling over atom and cell classes
    gives a stochastic fit objective; an independent fixed validation
    sample measures its error.
    \item \textbf{Constraint nullspace.} For constraint row $r$, compute
    the short vector
    $b_r^T=(C_nU)_{r,:}$ from one-leg atom and image sums. Stream the rows
    and accumulate
    \begin{equation}
        G_C
        =
        \sum_r b_rb_r^T
        =
        (C_nU)^T(C_nU).
        \label{eq:constraint-gram}
    \end{equation}
    $G_C$ has size $R_n\times R_n$. Its zero-eigenvalue eigenvectors form
    the nullspace basis $Z$ because
    $\operatorname{null}(G_C)=\operatorname{null}(C_nU)$.
\end{enumerate}

Permutation rows are absent in the symmetric-power fit because each
candidate is already permutation invariant. Translation covariance is
implemented by the common-origin sum or relative-image gauge. Only the
partition and explicit ASR image-sum rows need to be streamed into
\autoref{eq:constraint-gram}.

The memory held by the fitter is then
\begin{equation}
    O(R_nM+R_n^2)
    \label{eq:matrixfree-fit-memory}
\end{equation}
plus geometry-oracle and minibatch storage, instead of $O(M^n)$. The dense
and matrix-free implementations must agree on small cells before the latter
is used for production fits.

\section{Fine-q TD-SCHA averages from symmetric kernel factors}
\label{sec:fine-averages}

The low-rank kernel is useful only if its factors can be inserted directly
into the TD-SCHA stochastic averages. This section gives that construction.
The equal-factor, permutation-symmetric representation is assumed:
\begin{equation}
    K^{(n)}
    \simeq
    \sum_{\xi=1}^{R_n}c_\xi
    (A^{(\xi)})^{\otimes n}.
    \label{eq:fine-kernel-symmetric}
\end{equation}
The same geometry window $A^{(\xi)}$ acts on every tensor leg, while the
physical field, momentum, mode, and force/displacement role of each leg
remain different.

\subsection{Fine mesh, pair list, and windowed fields}
\label{sec:fine-fields}

Let $\mathcal Q_f$ be a uniform fine mesh containing $\Nf$ wavevectors.
For a perturbation at $\bm q_p$, define the fine two-phonon pair set
\begin{equation}
    \mathcal P_f(\bm q_p)
    =
    \{(\bm q_1,\bm q_2)\in\mathcal Q_f^2:
    \bm q_1+\bm q_2=\bm q_p+\bm G\}.
    \label{eq:fine-pair-set}
\end{equation}
Only one member of each exchanged pair is stored. Its multiplicity is one
when $\bm q_1=\bm q_2$ and two otherwise. A hash or integer mesh-index map
constructs \autoref{eq:fine-pair-set} in $O(\Nf)$ time.

Let $m=1,\ldots,N_{\mathrm{conf}}$ label the coarse-supercell stochastic
configurations and let $\rho_m$ be their normalized statistical weights.
The real-space displacement and force-residual fields are
$u_a^\alpha(\bm R;m)$ and $\delta f_a^\alpha(\bm R;m)$. For one kernel
factor $A^{(\xi)}$, define its fine-q Cartesian fields
\begin{align}
    \widetilde u_{a\alpha}^{\xi}(\bm q;m)
    &=
    \frac{1}{\sqrt{\Nc}}
    \sum_{\bm R,\bm p}
    A_{a\alpha\bm R}^{(\xi)}(\bm p)
    e^{-2\pi i\bm q\cdot[\bm R+\bm L(\bm p)]}
    u_a^\alpha(\bm R;m),
    \label{eq:fine-windowed-u}\\
    \widetilde f_{a\alpha}^{\xi}(\bm q;m)
    &=
    \frac{1}{\sqrt{\Nc}}
    \sum_{\bm R,\bm p}
    A_{a\alpha\bm R}^{(\xi)}(\bm p)
    e^{-2\pi i\bm q\cdot[\bm R+\bm L(\bm p)]}
    \delta f_a^\alpha(\bm R;m).
    \label{eq:fine-windowed-f}
\end{align}
The sums use the coarse data repeatedly with the image-dependent phase and
window. They do not require a fine-supercell ensemble.

The interpolated SSCHA dynamical matrix supplies the fine-q polarization
vectors $e_{a\alpha,\nu}(\bm q)$ and frequencies $\omega_\nu(\bm q)$. The
mode-space fields passed to the TD-SCHA kernel are
\begin{align}
    x_{\nu}^{\xi}(\bm q;m)
    &=
    \sum_{a\alpha}
    e_{a\alpha,\nu}^{*}(\bm q)
    \sqrt{M_a}\,
    \widetilde u_{a\alpha}^{\xi}(\bm q;m),
    \label{eq:fine-mode-x}\\
    y_{\nu}^{\xi}(\bm q;m)
    &=
    \sum_{a\alpha}
    e_{a\alpha,\nu}^{*}(\bm q)
    \frac{\widetilde f_{a\alpha}^{\xi}(\bm q;m)}{\sqrt{M_a}}.
    \label{eq:fine-mode-y}
\end{align}
$M_a$ is the atomic mass. One factor table therefore produces one pair of
fine-q field arrays $(x^\xi,y^\xi)$, evaluated at every momentum needed by
the pair list.

\subsection{General low-rank stochastic average}
\label{sec:general-fine-average}

Let $g_s^\xi(\bm q_s;m)$ denote either an $x^\xi$ field, a $y^\xi$ field,
or the same field after multiplication by a diagonal covariance factor.
Multilinearity and \autoref{eq:fine-kernel-symmetric} give the centered
rank-$n$ average
\begin{align}
    &\left\langle
    g_1(\bm q_1)\cdots g_n(\bm q_n)
    \right\rangle_K
    \nonumber\\
    &\quad =
    s_n
    \sum_{\xi=1}^{R_n}c_\xi
    \sum_{m=1}^{N_{\mathrm{conf}}}\rho_m
    \prod_{s=1}^{n}g_s^\xi(\bm q_s;m),
    \qquad
    \sum_{s=1}^{n}\bm q_s=\bm G,
    \label{eq:general-low-rank-average}
\end{align}
with the mesh-normalization factor
\begin{equation}
    s_n=
    \left(\frac{\Nc}{\Nf}\right)^{(n-2)/2},
    \qquad
    s_3=\sqrt{\frac{\Nc}{\Nf}},
    \qquad
    s_4=\frac{\Nc}{\Nf}.
    \label{eq:fine-vertex-scaling}
\end{equation}
The factor $s_n$ converts the normalization of a vertex estimated from the
coarse ensemble into the normalization of an $\Nf$-point fine-mesh vertex.
It multiplies the stochastic contraction once; it is not part of the
geometry-window normalization.

Equation~\eqref{eq:general-low-rank-average} is the central runtime result.
It replaces one contraction with a stored rank-$n$ tensor by a sum over
$R_n$ products of one-leg fine-q fields. The full kernel and the high-order
force-constant tensor are both absent.

\subsection{D3 contribution to the Lanczos action}
\label{sec:fine-d3-action}

Define the covariance-inverse-weighted displacement field
\begin{equation}
    r^\xi(\bm q;m)
    =
    f_Y(\bm q)\odot x^\xi(\bm q;m),
    \qquad
    f_Y=\frac{2\omega}{1+2n_B},
    \label{eq:fine-r-field}
\end{equation}
where $\odot$ is a mode-by-mode product and $n_B$ is the Bose
occupation evaluated at the fine-q frequency and temperature. Let
$R^{(1)}(\bm q_p)$ be the
one-phonon component of the Lanczos input vector. For each factor and
configuration, form the two scalar contractions
\begin{align}
    w_R^{\xi m}
    &=
    \frac{\rho_m}{3}
    \sum_\nu
    \overline{r_\nu^\xi(\bm q_p;m)}
    R_\nu^{(1)}(\bm q_p),
    \label{eq:d3-weight-r}\\
    w_F^{\xi m}
    &=
    \frac{\rho_m}{3}
    \sum_\nu
    R_\nu^{(1)}(\bm q_p)
    \overline{y_\nu^\xi(\bm q_p;m)}.
    \label{eq:d3-weight-f}
\end{align}
For every fine pair $(\bm q_1,\bm q_2)\in\mathcal P_f(\bm q_p)$, the D3
part of the two-phonon output has the schematic matrix form
\begin{align}
    d_2v_3(\bm q_1,\bm q_2)
    ={}&-s_3
    \sum_{\xi,m}c_\xi
    \bigl\{
    w_R^{\xi m}
    [r^\xi(\bm q_1;m)y^\xi(\bm q_2;m)^T
    +y^\xi(\bm q_1;m)r^\xi(\bm q_2;m)^T]
    \nonumber\\
    &\hspace{31mm}
    +w_F^{\xi m}
    r^\xi(\bm q_1;m)r^\xi(\bm q_2;m)^T
    \bigr\}.
    \label{eq:fine-d3-action}
\end{align}
The three possible positions of the force field are already present in
\autoref{eq:fine-d3-action}. Because all legs use the same $A^{(\xi)}$,
their sum is permutation symmetric without separate window permutations.
The adjoint map from the two-phonon input blocks to the one-phonon output
uses the same field products and coefficients $c_\xi$. Reusing exactly the
same factors on the two maps preserves the Hermiticity of the Lanczos
operator.

\subsection{D4 contribution without a four-rank tensor}
\label{sec:fine-d4-action}

Let $\alpha^{(1)}(\bm q_a,\bm q_b)$ be a two-phonon input block, with
$(\bm q_a,\bm q_b)\in\mathcal P_f(\bm q_p)$. For one factor and
configuration, first contract the two internal legs over the complete fine
pair list. Define
\begin{align}
    S_{xx}^{\xi m}
    &=
    \sum_{(\bm q_a,\bm q_b)\in\mathcal P_f}
    \mu_{ab}
    x^\xi(\bm q_a;m)^\dagger
    \alpha^{(1)}(\bm q_a,\bm q_b)
    \overline{x^\xi(\bm q_b;m)},
    \label{eq:d4-internal-xx}\\
    b^\xi(\bm q_a;m)
    &=
    \alpha^{(1)}(\bm q_a,\bm q_b)
    \overline{x^\xi(\bm q_b;m)},
    \label{eq:d4-buffer}\\
    S_{xy}^{\xi m}
    &=
    \sum_{\bm q_a,\nu}
    \overline{b_\nu^\xi(\bm q_a;m)}
    f_\psi(\bm q_a,\nu)
    y_\nu^\xi(\bm q_a;m),
    \qquad
    f_\psi=\frac{1+2n_B}{2\omega}.
    \label{eq:d4-internal-xy}
\end{align}
$\mu_{ab}$ is the stored-pair multiplicity. The buffer in
\autoref{eq:d4-buffer} is also accumulated at $\bm q_b$ with the transpose
of $\alpha^{(1)}$ when the two momenta differ.

The external two-phonon block is then
\begin{align}
    d_2v_4(\bm q_1,\bm q_2)
    ={}&-s_4
    \sum_{\xi,m}c_\xi\rho_m
    \biggl\{
    \frac{S_{xx}^{\xi m}}{8}
    [r^\xi(\bm q_1;m)y^\xi(\bm q_2;m)^T
    +y^\xi(\bm q_1;m)r^\xi(\bm q_2;m)^T]
    \nonumber\\
    &\hspace{27mm}
    +\frac{S_{xy}^{\xi m}}{4}
    r^\xi(\bm q_1;m)r^\xi(\bm q_2;m)^T
    \biggr\}.
    \label{eq:fine-d4-action}
\end{align}
The first line places the force field on either external leg. The second
line comes from placing it on either internal leg through
$S_{xy}^{\xi m}$. These are the four D4 force-position channels already
implemented in the q-space kernel. Since the same factor $A^{(\xi)}$ is
used in all four positions, one factor pass is permutation symmetric and
the current reused $w,v$ data path is sufficient.

The expensive internal pair sum is performed once for each
$(\xi,m)$ and reused for every external pair. Therefore one D4 factor pass
remains $O(\Nf n_b^2)$ rather than $O(\Nf^2 n_b^4)$. At no point is a
four-index vertex over atom, image, or mode labels stored.

\subsection{Prefiltering and fine-side analytic factors}
\label{sec:fine-prefilter}

The raw displacement correlations contain the long-ranged SSCHA covariance
$f_\psi$. The repository tests show that windowing those raw correlations
can bias the vertex. The production path must therefore keep the existing
prefiltering order:
\begin{enumerate}
    \item apply $f_Y=f_\psi^{-1}$ exactly on the commensurate coarse mesh;
    \item transform the filtered configurations back to coarse real space;
    \item apply the low-rank image windows and evaluate
    \autoref{eq:fine-windowed-u} on the fine mesh;
    \item obtain fine-q frequencies, polarizations, $f_Y$, and $f_\psi$
    from the interpolated SSCHA dynamical matrix;
    \item fold the required fine-side $f_\psi$ factors into the
    $\alpha^{(1)}$ blocks.
\end{enumerate}
This order makes the windowed stochastic correlation decay with the
force-constant range rather than the phonon-propagator range. It is an
exact rewriting at commensurate q points and must be applied identically to
every low-rank factor.


\section{Implementation strategy and cost}
\label{sec:implementation}

\subsection{Primary data representation}
\label{sec:code-changes}

The production implementation should support the fully symmetric form
first. For order $n=3$ or $4$, the persistent fitting result is
\begin{equation}
    \mathcal F_n
    =
    \{c_\xi,A^{(\xi)}:\xi=1,\ldots,R_n\}.
    \label{eq:factor-data-product}
\end{equation}
Each $c_\xi$ is a signed scalar and each $A^{(\xi)}$ is a one-leg table
indexed by primitive atom, Cartesian component when required, folded
coarse cell, and supported image. Geometry-only fits normally share the
same factor among Cartesian components, so that index should be omitted
unless a later weighted fit requires it. The file containing
$\mathcal F_n$ must also record the lattice, atomic positions, supercell,
image support, target metric, ASR convention, fit tolerance, and a format
version. A hash of these quantities prevents factors from being reused for
the wrong structure.

The following objects are explicitly forbidden in the production path:
the dense $T^{(4)}$, $K^{(4)\star}$, $K^{(4)}$, $C_4$, and
$\Phi^{(4)}$. They may exist only behind a small-system validation flag.

\subsection{Setup pipeline}
\label{sec:setup-pipeline}

The implementation can be divided into six stages.

\paragraph{1. Geometry and constrained factor fit.}
Build neighbor lists and the complete-graph image oracle from the primitive
structure. For a small validation system, construct the dense target and
compare it with the matrix-free oracle entry by entry. For production, fit
the symmetric factors with the operations in \autoref{sec:matrix-free}:
stream target tuples, build one-leg Gram matrices, accumulate the constraint
Gram matrix, and update the factors without constructing a four-leg array.
Increase $R_n$ adaptively until the kernel error reaches the stochastic
noise target or the runtime-rank budget is exceeded.

\paragraph{2. Fine mesh and harmonic quantities.}
Construct $\mathcal Q_f$ and the integer pair map
$\mathcal P_f(\bm q_p)$. Interpolate the SSCHA dynamical matrix to every
fine q point, reuse the original matrix at commensurate points, impose the
rank-two ASR, and diagonalize it to obtain
$\omega_\nu(\bm q)$ and $e_{a\alpha,\nu}(\bm q)$. Store q points by integer
mesh coordinates so momentum partners are found by modular arithmetic
rather than floating-point searches.

\paragraph{3. Coarse-side prefilter.}
For every stochastic configuration, transform the displacement field to
the coarse commensurate mesh, apply the exact mode-diagonal $f_Y$, and
transform it back to the coarse real-space representation. Keep the force
residual in the convention already used by the q-space TD-SCHA kernel.
This stage is independent of the fitted rank and is performed once per
configuration batch.

\paragraph{4. Windowed fine-q fields.}
For factor $\xi$, evaluate \autoref{eq:fine-windowed-u} and
\autoref{eq:fine-windowed-f}, then rotate them with the fine-q
polarizations using \autoref{eq:fine-mode-x} and
\autoref{eq:fine-mode-y}. The phase
$e^{-2\pi i\bm q\cdot[\bm R+\bm L(\bm p)]}$ depends on q, folded cell,
and image but not on the stochastic configuration. Precompute it in q
chunks or evaluate it by a batched non-uniform discrete Fourier transform.
Do not assemble a dense matrix with simultaneous factor, q, atom, cell,
and image indices when that matrix exceeds the field-memory budget.

\paragraph{5. Anharmonic Lanczos application.}
Stream the factor list. For each $\xi$, pass the same field set
$(x^\xi,y^\xi)$ to every logical leg and evaluate the existing D3 and D4
force-position channels. Multiply the returned action by $c_\xi s_3$ or
$c_\xi s_4$ and accumulate it immediately. The D3 forward and adjoint maps
must use the identical factors and coefficients. The D4 internal scalars
in \autoref{eq:d4-internal-xx} and \autoref{eq:d4-internal-xy} are computed
once per factor, configuration, and symmetry, then reused over all external
pairs.

\paragraph{6. Parallel reduction and persistence.}
Distribute configurations and crystal symmetries as in the current code.
The factor index may be distributed as a second level when $R_n$ is large,
provided the partial one-phonon and two-phonon outputs are reduced only
after multiplication by $c_\xi$. Persist the factors and diagnostics, not
the fine-q fields. Fine fields depend on the ensemble batch and should be
recomputed or cached only within a Lanczos run.

\subsection{Chunking policy}
\label{sec:chunking}

Three independent chunk sizes should be exposed:
\begin{align}
    B_c&=\text{configurations per batch},\nonumber\\
    B_q&=\text{fine q points per batch},\nonumber\\
    B_R&=\text{kernel factors per batch}.
    \label{eq:chunk-sizes}
\end{align}
The default should keep one or a few factors in memory
($B_R\ll R_n$), because equal-factor terms do not need simultaneous access
to other ranks. $B_q$ must include complete momentum partners or use a
pair-index indirection into the current q chunk. $B_c$ follows the existing
MPI configuration distribution. A memory estimator should choose safe
defaults before allocating the field arrays.

\subsection{Validation gates}
\label{sec:implementation-gates}

Implementation should proceed only through the following gates:
\begin{enumerate}
    \item dense and matrix-free target entries agree on a small cell;
    \item streamed partition and ASR residuals agree with dense
    $C_nk-d_n$ and reach numerical precision;
    \item at all commensurate q points, each assembled D3 and D4 action
    equals the original coarse operator configuration by configuration;
    \item the D3 forward and adjoint maps satisfy the Lanczos Hermiticity
    check;
    \item factor-rank convergence is demonstrated on a toy system with a
    direct fine-supercell reference;
    \item the SnTe off-grid D3 and D4 benchmarks pass before the mode is
    exposed as a production default.
\end{enumerate}

\subsection{Runtime and memory}
\label{sec:runtime}

This subsection gives costs for the primary equal-factor implementation.
Define
\begin{equation}
    D_c=3n_{\mathrm{at}}\Nc,
    \qquad
    J=N_{\mathrm{img}},
    \qquad
    M=D_cJ,
    \qquad
    R=R_3+R_4.
    \label{eq:cost-symbols}
\end{equation}
$D_c$ is the number of Cartesian degrees of freedom in the coarse
supercell, $J$ is the supported image count per leg, $M$ is the maximum
one-leg factor size, and $R$ is the total number of symmetric D3 and D4
factors. Geometry-only factors that omit the Cartesian index are smaller by
a factor of three.

\paragraph{Persistent factor storage.}
Storing the signed factor lists requires
\begin{equation}
    O(RM)
    \label{eq:factor-storage}
\end{equation}
real numbers. This is the principal advantage over $O(M^4)$ D4 storage.
The factors are independent of $\Nf$ and of the number of stochastic
configurations.

\paragraph{Fine harmonic data.}
Frequencies require $O(\Nf n_b)$ storage and polarization matrices require
$O(\Nf n_b^2)$. For a large primitive cell and dense q mesh, polarization
storage can exceed the low-rank field storage. The q points should therefore
be distributed across processes or stored in chunked on-disk arrays. The
pair list itself is $O(\Nf)$, while the two-phonon Lanczos vector contains
$O(\Nf n_b^2)$ complex numbers; this cost already exists in the direct
fine-q formulation and is independent of the centering fit.

\paragraph{Field construction.}
Building all windowed Cartesian and mode fields directly has the
upper-bound cost
\begin{equation}
    O\left[
    R N_{\mathrm{conf}}\Nf(M+n_b^2)
    \right].
    \label{eq:field-build-cost}
\end{equation}
The $M$ term is the image-weighted real-space sum and the $n_b^2$ term is
the fine-q polarization rotation. Batched Fourier transforms, factor
sparsity, shared image phases, and crystal symmetries reduce the prefactor.
The field builder should report the achieved cost separately from the
Lanczos contraction because it may dominate when $n_b$ is large.

\paragraph{Lanczos application.}
Once the fields are available, one anharmonic application scales as
\begin{equation}
    O\left[
    (R_3+R_4)
    N_{\mathrm{conf}}N_{\mathrm{sym}}\Nf n_b^2
    \right].
    \label{eq:runtime-cost}
\end{equation}
For D4, the internal pair reduction and the external block update are both
linear in the number of fine pairs. The internal scalar is reused across
all external pairs, as shown in \autoref{eq:fine-d4-action}. The symmetric
factor has no $n!$ permutation multiplier. All force-position channels are
accumulated inside the same factor pass.

\paragraph{In-memory field cache.}
If all q points for a batch are cached, one displacement or force family
requires
\begin{equation}
    O(B_RB_c\Nf n_b)
    \label{eq:full-field-cache}
\end{equation}
complex numbers. With q chunking this becomes
\begin{equation}
    O(B_RB_cB_qn_b).
    \label{eq:chunked-field-cache}
\end{equation}
The D4 operation can use two streamed sweeps: the first accumulates
$S_{xx}^{\xi m}$ and $S_{xy}^{\xi m}$ over internal-pair chunks, and the
second forms the external blocks. This avoids retaining all factor fields
at once. If q fields are not cached between the sweeps, they are recomputed;
the memory estimator must expose this time-memory tradeoff.

\paragraph{Scaling limits.}
The low-rank method removes the exponential storage in tensor order, but it
does not remove the intrinsic $n_b^2\Nf$ size of the two-phonon problem.
For very large primitive cells, the band-pair dimension is the limiting
cost. For very dense meshes, q-distributed pair blocks and field streaming
are required. The implementation should refuse an allocation whose
predicted peak memory exceeds a user-provided fraction of available memory
and should print the factor, polarization, pair-block, and field-cache
contributions separately.
\subsection{Setup cost and fallback policy}
\label{sec:setup-cost}

The setup must be budgeted by oracle evaluations and factor operations, not
by dense-kernel dimensions. Let $N_{\mathrm{oracle}}$ be the number of
folded atom-cell tuples evaluated in one optimization sweep and let
$N_{\mathrm{it}}$ be the number of factor sweeps. For symmetric rank $R_n$,
the main setup operations are:
\begin{itemize}
    \item $O(R_n^2M)$ work for one-leg factor Gram matrices;
    \item streamed evaluation of $N_{\mathrm{oracle}}$ geometry targets and
    their contractions with $R_n$ symmetric powers;
    \item accumulation and diagonalization of the
    $R_n\times R_n$ constraint Gram matrix in
    \autoref{eq:constraint-gram};
    \item factor-gradient and line-search work over $O(R_nM)$ parameters.
\end{itemize}
The total is proportional to $N_{\mathrm{it}}$, but never requires an
$M^4$ allocation. A naive geometry-oracle call scans
$N_{\mathrm{img}}^{n-1}$ relative-image combinations. Complete-graph bounds,
nearest-image candidate lists, exact-tie detection, and neighbor lists must
prune that search before rank-four production fitting.

For small cells, the dense target remains the reference implementation. A
configuration option should set strict limits on the maximum dense entry
count and memory. Crossing either limit must select the matrix-free path
automatically rather than attempting the allocation.

The fitted factors are geometry data and should be cached. A production run
loads $\mathcal F_3$ and $\mathcal F_4$, verifies their structure hash and
constraint diagnostics, and proceeds directly to the fine-field setup. A
failed or interrupted fit should checkpoint the factor values, optimizer
state, random-sampling seed, and accumulated validation sample so it can be
resumed reproducibly.

The fallback order is:
\begin{enumerate}
    \item fit the permutation-symmetric equal-factor model;
    \item accept it only if the geometry error, ASR, commensurate identity,
    and estimated runtime-product budget all pass;
    \item if its required rank is too large, fit the distinct-factor model
    on the same oracle samples and compare measured ordered products,
    distinct fields, and wall time;
    \item introduce the four-slot Julia interface only if the distinct-factor
    model is demonstrably cheaper at the same physical accuracy;
    \item use an explicit centered $\Phi^{(4)}$ only as a small-system
    validation oracle.
\end{enumerate}
This policy keeps the already permutation-symmetric case as the production
design and prevents a speculative general interface from becoming the
default implementation cost.

\section{Validation and known pitfalls}
\label{sec:validation}

The repository history supplies several negative controls that must remain
visible in the implementation:
\begin{enumerate}
    \item a plain off-grid window is exact at commensurate q but can be badly
    biased between those points;
    \item independently centered rank-three factors can violate ASR, while
    projecting windowed stochastic fields afterward changes the fitted
    kernel and can bias the response;
    \item shifting a compact window is not equivalent to shifting the
    stochastic data when support wraps around the supercell;
    \item a cell-index target fails for basis-dependent ties in an
    $L=2$ real-material cell;
    \item a force-star D4 target can pass channel-sum and commensurate
    identities while representing the wrong four-leg geometry;
    \item the present D4 reference and pinned-force-leg modes are diagnostic
    baselines, not a validated solution.
\end{enumerate}

The validation ladder for a new fit is:
\begin{enumerate}
    \item geometry tests for exact ties, translation covariance, triangle or
    complete-graph minima, and permutation symmetry;
    \item machine-precision partition and ASR residuals for every leg;
    \item configuration-by-configuration equality to the current operator
    on the complete commensurate coarse mesh;
    \item Hermiticity and adjointness of the assembled Lanczos action;
    \item acoustic-leg vertices that vanish at $\Gamma$ and have the correct
    small-q scaling without relying on a posteriori projection;
    \item rank convergence against a toy model where a direct fine-supercell
    rank-three or rank-four calculation is affordable;
    \item the SnTe $2^3\rightarrow4^3$ off-grid spectral and static-response
    benchmarks, with the old plain, reference, and force-star modes retained
    as negative controls.
\end{enumerate}

The fit must stop at the geometry gate if the projection distance is large,
and at the rank gate if the projected target needs too many runtime terms.
These failures identify insufficient support or an unsuitable target; they
must not be hidden by relaxing commensurate identity or ASR.

\section{Conclusions}
\label{sec:conclusions}

The low-rank window formalism is a viable tensor-free interpolation strategy
provided that the fit is performed at the kernel level. Partition of unity
preserves the commensurate tensor, the explicit image-sum rows preserve ASR,
and equal-factor powers impose permutation symmetry term by term. These
properties constrain the assembled kernel correction and remain exact
during the locality fit.

For rank three, the geometry target is a triangle. For rank four, it is a
genuine four-leg complete graph with six pair distances. The symmetric-power
representation can use the current D3 force channels and reused D4 $w,v$
path. A distinct-factor four-slot implementation is only a measured
fallback when it provides lower total runtime at the same constrained
kernel error.

The fine-q TD-SCHA action is obtained by transforming each one-leg factor
into windowed displacement and force fields, then inserting those fields
into the existing three- and four-field stochastic averages. The factors
$s_3=\sqrt{\Nc/\Nf}$ and $s_4=\Nc/\Nf$ supply the fine-mesh vertex
normalization. Internal D4 pair contractions are reduced to scalars before
the external blocks are formed, keeping every factor pass linear in
$\Nf$.

A dense projected D4 kernel is a validation object only. Production fitting
uses the geometry oracle, one-leg Gram matrices, streamed ASR and partition
rows, and an $R_n\times R_n$ constraint Gram matrix. Its storage is
$O(R_nM+R_n^2)$ rather than $O(M^4)$. Fine-mesh fields, polarizations, and
two-phonon blocks are independently chunked or distributed because their
cost grows with $\Nf$ and $n_b^2$. The decisive gates are therefore
constrained factor rank, matrix-free geometry error, fine-field memory, and
measured Lanczos wall time.

\end{document}
